\section{線形倒立振子モデル（LIPM）の概要}

線形倒立振子モデル（LIPM）は，二足歩行ロボットの歩行制御において広く用いられている簡略化モデルである．LIPMでは，ロボットの全質量が一点のCoMに集中していると仮定し，そのCoMが一定の高さで運動する倒立振子としてモデル化される．また，角運動量の影響を無視することで，CoMの水平方向運動とZMPの関係を線形な形で記述できる．

このモデルにより，ZMPを入力としてCoM運動を解析・設計することが可能となり，歩行生成や安定化制御が大幅に簡略化される．一方で，CoM高さが一定であるという仮定は，床が剛体であることを前提としており，床の変形や沈み込みが生じる状況では必ずしも成立しない．

LIPMは実用上非常に有用なモデルであるが，その仮定と適用範囲を理解したうえで使用する必要がある．本研究では，この仮定が破綻する状況として，軟弱地面における歩行を対象とする．